\documentclass[addpoints]{exam}
% \documentclass[addpoints,12pt]{exam}

\usepackage[slovene]{babel}
\usepackage{amsfonts}
\usepackage{amssymb}
\usepackage{amsmath}
\usepackage{float}
\usepackage{graphicx}
\usepackage{enumitem}
\usepackage{color}
\usepackage[gen]{eurosym}
\usepackage{newunicodechar}
\newunicodechar{€}{\euro}


%Komentar naslednje vrstice glede na verzijo testa -- rešitve ali prostor za odgovore:
% \printanswers

% \linespread{2}


\pointsinrightmargin
\marginbonuspointname{}


\renewcommand{\solutiontitle}{\noindent\textbf{Rešitve in točkovanje:}\par\noindent}

\colorgrids
\definecolor{GridColor}{gray}{0.7}
\setlength{\gridsize}{7mm}

\extrawidth{5mm}
\setlength{\rightpointsmargin}{1.5cm}

\begin{document}

\runningfooter{}{}{}

\begin{center}
    \fbox{\fbox{\parbox{15cm}{\centering
    Četrto pisno ocenjevanje znanja \\ 1. letnik -- test A \\ 16. april 2025 \\ (Realna števila)}}}
\end{center}

\vspace{0.1in}
\makebox[\textwidth]{Ime in priimek, oddelek:\enspace\hrulefill}


\begin{center}
    \chqword{Naloga}
    \chpword{Možne točke}
    \chbpword{Dodatne točke}
    \chsword{Dosežene točke}
    \chtword{Skupaj}  
    % \combinedpointtable[h][questions]
    \combinedgradetable[h][questions]

\end{center}

\vspace{0.3in}


\begin{questions}

    % 1. vprašanje

    \question[6]
    Natančno izračunajte vrednost izraza.
    $$ \dfrac{\sqrt{5}}{\sqrt{5}+1}-\dfrac{5-3\sqrt{5}}{4\sqrt{5}}+\left(3+\sqrt{5}\right)^2-\sqrt{125} $$

    \begin{solution}[6cm]
        ????
    \end{solution}
    
    
    \newpage

    % 2. vprašanje
    \question[4]
    Rešite enačbo.
    $$ \dfrac{2}{x-3}-\dfrac{x}{x-1}=\dfrac{2x}{x^2-4x+3} $$
    
    
    \begin{solution}[9cm]
        ???
    \end{solution}

        


    % 3. vprašanje
    \question[5]
    Rešite sistem neenačb, rešitev grafično upodobite in zapišite z intervalom.
    $$ \left(x+1>3\right)\land\left(2x\leq 3(x-1)\right) $$
    
    
    \begin{solution}[8cm]
        ???
    \end{solution}


    \newpage 


    % 4. vprašanje
    \question[4]
    Rešite sistem enačb.
    $$ \begin{aligned}
        y-x&=3 \\ x-y&=5
    \end{aligned} $$
    
    
    \begin{solution}[9cm]
        ???
    \end{solution}



    % 5. vprašanje
    \question[4]
    Prva cev sama napolni bazen v $6$ urah, druga sama pa v $9$ urah. Drugo cev odpremo eno uro kasneje kot prvo.
    V kolikšnem času bo bazen poln?
    
    
    \begin{solution}[8cm]
        ???
    \end{solution}




    \newpage 

    
    % 6. vprašanje
    \question[4]
    Natančno izračunajte.
    $$ \left\lvert \lvert\sqrt{2}+4\rvert-\sqrt{2}\right\rvert + \left\lvert \lvert\sqrt{2}-4\rvert-\sqrt{2}\right\rvert $$
    
    
    
    \begin{solution}[10cm]
        ???
    \end{solution}




    % 7. vprašanje
    \question[5]
    Rešite enačbo in opišite njen geometrijski pomen ter rešitev grafično prikažite.
    $$  \left\lvert x+3\right\rvert=2 $$
    
    
    \begin{solution}[8cm]
        ???
    \end{solution}


    \newpage

    \makebox[\textwidth]{Ime in priimek, oddelek:\enspace\hrulefill}

    ~\\



    % 8. vprašanje
    \question[5]
    Izračunajte vrednosti izraza glede na vrednosti spremenljivke $x$.
    $$\left\lvert x+8 \right\rvert + \left\lvert x-2 \right\rvert $$
    
    
    \begin{solution}[10cm]
        ???
    \end{solution}



    % 9. vprašanje
    \question[4]
    V posodi imamo $220~g$ $95~\%$ raztopine kisline. Koliko gramov $60~\%$ raztopine kisline moramo dodati, da bomo dobili $80~\%$ raztopino kisline?
    
    
    \begin{solution}[7cm]
        ???
    \end{solution}



    \newpage

    % 10. vprašanje
    \question[4]
    Izlet v Grčijo se je najprej podražil za $5~\%$, nato pa še za $10~\%$. 
    Koliko je stal na začetku, če zdaj stane $2310~€$?
    
    
    \begin{solution}[3cm]
        ???
    \end{solution}

     \newpage



    % 3. vprašanje
    % \question[4]
    % Produkt dveh števil je $192$, njun največji skupni delitelj pa $4$. Poiščite števili.
    
    % \begin{solution}[7cm]
    %     \begin{align*}
    %         1-(1+x^8)(1+x^4)(1+x^2)(1+x)(1-x)&=1-(1+x^8)(1+x^4)(1+x^2)(1-x^2)= \\
    %             &=1-(1+x^8)(1+x^4)(1-x^4)= \\
    %             &=1-(1+x^8)(1-x^8)= \\
    %             &=1-(1-x^{16})= \\
    %             &=1-1+x^{16}= \\
    %             &=x^{16}
    %     \end{align*}
    %     Za pravilno uporabo pravila razlike kvadratov $1$ točka, za pravilno upoštevan negativen predznak $1$ točka, končen rezultat $1$ točka.
    % \end{solution}


    % \begin{description}
    %     \item[a] $(x-a)^2+(y-b)^2-c^2=0$ 
    %     \item[b] $d^2(x-a)^2+f^2(y-b)^2-(df)^2=0$
    %     \item[c] $S(a,b)$
    %     \item[d] $e^2=d^2-f^2;\ d>f$
    % \end{description}
    % Krožnica: \fillin[a c][2cm]
    % \\ Elipsa:  \fillin[b c d][2cm]




     \newpage

    %  \noindent Ime in priimek: \underline{~~~~~~~~~~~~~~~~~~~~~~~~~~~~~~~~~~~~~~~~~~~~~~~~}

    %  ~\\

    %dodatno vprašanje
    \bonusquestion[3]
    \textit{Dodatna naloga:} Obravnavajte enačbo. $$a(x-a)=2(x-1)-2$$

    \begin{solution}[10cm]
        ????
    \end{solution}




    % \newpage
        

    % dodatno vprašanje
    \bonusquestion[2]
    \textit{Dodatna naloga:} Natančno izračunajte. $$ \left(\sqrt{3}+\sqrt{5}\right)\sqrt{8-2\sqrt{15}} $$

    \begin{solution}[7cm]
        ????
    \end{solution}




\end{questions}



\end{document}